\documentclass[10pt, landscape, twocolumn]{article}

% --- Librerías Requeridas ---
\usepackage[margin=0.5in, columnsep=0.4in]{geometry}
\usepackage{amsmath, amssymb, amsthm}
\usepackage[table]{xcolor}
\usepackage{tabularx}
\usepackage{tikz}
\usetikzlibrary{arrows.meta, positioning, calc, shapes.geometric, fit, backgrounds}
\usepackage{enumitem}
\usepackage{mathpazo} % Tipografía Palatino
\usepackage{booktabs}

% --- Configuración de Formato ---
\pagestyle{empty} % Sin números de página
\setlength{\parindent}{0pt}
\setlength{\parskip}{2pt}
\raggedbottom

% Ajuste de espacios en ecuaciones
\AtBeginDocument{
  \setlength{\abovedisplayskip}{3pt}
  \setlength{\belowdisplayskip}{3pt}
  \setlength{\abovedisplayshortskip}{1pt}
  \setlength{\belowdisplayshortskip}{1pt}
}

% --- Definición de Colores ---
\definecolor{themeblue}{RGB}{0, 90, 160}
\definecolor{themegray}{RGB}{80, 80, 80}
\definecolor{setblue}{RGB}{240, 244, 255}
\definecolor{setred}{RGB}{255, 240, 240}
\definecolor{setborderblue}{RGB}{0, 60, 140}
\definecolor{setborderred}{RGB}{150, 0, 0}

% --- Comandos Personalizados para Títulos ---
\newcommand{\cheatsheetsection}[1]{
    \vspace{0.6em}
    {\color{themeblue}\large\bfseries #1} \par
    \vspace{-0.5em}
    {\color{themeblue}\rule{\linewidth}{1pt}} \par
    \vspace{0.2em}
}

\newcommand{\cheatsheetsubsection}[1]{
    \vspace{0.4em}
    {\color{themeblue}\normalsize\bfseries #1} \par
    \vspace{0.1em}
}

% --- Inicio del Documento ---
\begin{document}

% --- Cabecera ---
\begin{center}
    \vspace{-1em}
    {\color{themeblue}\huge\bfseries Título del Apunte o Tema}\\[0.1em]
    {\color{themegray}\small Subtítulo descriptivo o síntesis de los temas cubiertos}\\[0.4em]
    {\color{themegray}\footnotesize \textbf{Por Roberto León Landeros}}
\end{center}
\vspace{-0.8em}

% --- Sección 1 ---
\cheatsheetsection{1. Fundamentos y Definiciones Formales}
Texto directo y riguroso. Definir objetos formalmente con $:=$ (por ejemplo, sea $f : \mathbb{R}^n \to \mathbb{R}$ diferenciable tal que $\forall \mathbf{x} \in \Omega$, $\nabla f(\mathbf{x}) := (\partial_{x_1} f, \dots, \partial_{x_n} f)^T$).
\[ \int_{a}^{b} f(x) \, dx = F(b) - F(a) \]

\cheatsheetsubsection{1.1 Propiedades Estructurales}
\begin{itemize}[leftmargin=*, nosep]
    \item \textbf{Concepto Clave 1:} Explicación precisa y concisa desde primeros principios.
    \item \textbf{Concepto Clave 2:} Condición necesaria y suficiente ($\iff$) o implicación ($\implies$).
    \item \textbf{Caso Límite / Observación:} Nota técnica sobre estabilidad, complejidad o contraejemplos.
\end{itemize}

% --- Ejemplo de Tabla ---
\cheatsheetsection{2. Síntesis Comparativa}
\vspace{0.2em}
\renewcommand{\arraystretch}{1.2}
\begin{tabularx}{\linewidth}{|p{0.28\linewidth}|X|}
\hline
\rowcolor{setblue}
\textbf{Concepto / Operador} & \textbf{Definición Formal y Aplicación} \\ \hline
Elemento A & Descripción formal, complejidad o interpretación geométrica. \\ \hline
Elemento B & Restricciones, invariantes o casos especiales relevantes. \\ \hline
Fórmula C & $\nabla \cdot \mathbf{B} = 0$ --- Interpretación física o algebraica. \\ \hline
\end{tabularx}

\end{document}
